% Propietario conservado: /Users/ruben/Documents/ChatGPT/jueces y controles/REVISION_ARGUMENTAL_APERTURAS_CIERRES_20260915/IV/ES/sections/acoplamiento_recuperacion.tex
\section{Acoplamiento aritmético-geométrico y recuperación bilateral}
\label{iv:sec:acoplamiento-recuperacion}

El registro dodecafásico de la sección~\ref{iv:sec:k-registro} admite
lecturas aritméticas, incidenciales y operatorias. Denominamos
\emph{constante holográfica de acoplamiento aritmético-geométrico}
al registro generado \(K\), conservando el orden de sus doce ventanas,
el origen de fase y la orientación. Su lectura racional
\(\kappa=K_{\mathrm{int}}/(1000^{12}-1)\) es una codificación de ese
registro; no reemplaza sus tipos ni la historia enriquecida que lo produce.
La inversión de Hadamard recupera su coordenada firmada; la lectura
incidencial determina la tétrada de Witt. Las construcciones siguientes
explicitan otras dos funciones: identificar operadores sobre su portador
cíclico y conservar información entre dos realizaciones isométricas.
Son operaciones posteriores a la realización del continuo, no reglas
que introduzcan constantes convencionales en APP--TRIT--TPK.

\subsection{Codificación exacta y referencia cíclica}
\label{iv:subsec:acoplamiento-codigo}

\begin{proposition}[Recuperación del registro desde su lectura racional]
Fijadas base \(b=1000\), longitud doce, origen y orientación, la aplicación
\[
 (k_1,\ldots,k_{12})\longmapsto
 \kappa=\frac{\sum_{j=1}^{12}k_jb^{12-j}}{b^{12}-1},
 \qquad 0\leq k_j<b,
\]
es inyectiva. Desde \(\kappa\) se recupera el entero
\(n=(b^{12}-1)\kappa\) y, por divisiones euclídeas sucesivas, los doce
bloques, incluidos los ceros iniciales. Un error absoluto estrictamente
menor que \(1/[2(b^{12}-1)]\) permite la misma recuperación por redondeo.
\end{proposition}
\begin{proof}
La expansión entera de longitud fija en base \(b\) es única. Dos enteros
distintos están separados al menos por uno; sus lecturas racionales
están separadas al menos por \(1/(b^{12}-1)\). La cota estricta sitúa el
dato aproximado dentro de la celda de redondeo de un único entero.
\end{proof}

Consideremos ahora \(K=(234,543,140,729,659,824,621,58,914,794,146,601)\),
el desplazamiento cíclico \(S\) de sus posiciones y
\(O_K=(K,SK,\ldots,S^{11}K)\) sobre \(\mathbb R^{12}\).
El desplazamiento de posiciones no se identifica con toda la evolución
enriquecida ni, por su sola periodicidad, con una simetría excepcional.

\begin{theorem}[Invertibilidad de la referencia orbital]
\label{iv:thm:acoplamiento-orbita}
La matriz \(O_K\) es invertible y
\[
 589609\|c\|^2\leq\|O_Kc\|^2\leq6263^2\|c\|^2.
\]
Su órbita centrada tiene rango once.
\end{theorem}
\begin{proof}
El Gram \(O_K^*O_K\) es circulante. Su primera fila, obtenida mediante
los doce productos escalares \(\langle K,S^jK\rangle\), es
\[
 \begin{aligned}
 (&4251257,2999323,3163385,3287920,3216553,3344743,\\
  &2950064,3344743,3216553,3287920,3163385,2999323).
 \end{aligned}
\]
La evaluación en los doce caracteres del ciclo proporciona los valores
propios
\[
 \begin{gathered}
 6263^2,\quad697225,\quad
 589609\ (2),\quad1407529\ (2),\quad1053157\ (2),\\
 (1248025-345420\sqrt3)\ (2),\qquad
 (1248025+345420\sqrt3)\ (2),
 \end{gathered}
\]
donde los paréntesis indican multiplicidades. El menor valor es
\(589609\): para la única comparación no inmediata,
\(658416>345420\sqrt3\), como muestra el cuadrado de ambos lados.
La desigualdad triangular de la transformada de Fourier y \(K_j\geq0\)
dan el máximo \((\sum K_j)^2=6263^2\).
La positividad de todos los modos prueba la invertibilidad y las cotas.
Centrar el registro elimina únicamente el modo uniforme.
\end{proof}

\begin{corollary}[Identificación de un operador por sus respuestas]
\label{iv:cor:acoplamiento-identificacion}
Un operador lineal \(T:\mathbb R^{12}\to\mathbb R^{12}\) queda determinado por
\(Y=(TK,TSK,\ldots,TS^{11}K)\), mediante \(T=YO_K^{-1}\).
Si \(TS=ST\), basta la respuesta vectorial completa \(y=TK\), pues
\(T=O_yO_K^{-1}\). Para una perturbación \(E\) de \(Y\),
\(\|\widehat T-T\|\leq\|E\|/\sqrt{589609}\).
\end{corollary}
\begin{proof}
Se tiene \(Y=TO_K\). Bajo conmutación, \(TS^jK=S^jy\), luego \(Y=O_y\).
La cota se deduce de
\(\|O_K^{-1}\|=1/\sqrt{589609}\).
Las normas son la euclídea y su norma operatoria subordinada.
\end{proof}

La respuesta \(y\) contiene doce componentes. La afirmación no atribuye
a un único número escalar la recuperación de un operador arbitrario:
precisa qué observación y qué conmutación hacen posible la identificación.

\subsection{Conservación y restitución de dos canales}
\label{iv:subsec:acoplamiento-bilateral}

Sean \(H,H'\) espacios de Hilbert y
\(\mathcal B,\mathcal C:H\to H'\) isometrías lineales acotadas.
Para \(a>0\), definimos
\[
 p=a^2+a+1,\qquad N=(2a+1)p=(a+1)^3+a^3,\qquad
 Q_-=a\mathcal B-\mathcal C,\quad Q_+=(a+1)\mathcal B+\mathcal C.
\]

\begin{theorem}[Balance bilateral y recuperación]
\label{iv:thm:acoplamiento-balance}
Se cumplen
\[
 (a+1)Q_-^*Q_-+aQ_+^*Q_+=NI_H,\qquad
 W_av=\frac1{\sqrt N}
 \begin{pmatrix}\sqrt{a+1}\,Q_-v\\\sqrt a\,Q_+v\end{pmatrix},
 \qquad W_a^*W_a=I_H.
\]
Los datos \(x=Q_-v\), \(y=Q_+v\) determinan
\[
 \mathcal Bv=\frac{x+y}{2a+1},\qquad
 \mathcal Cv=\frac{ay-(a+1)x}{2a+1},\qquad
 v=\mathcal B^*\frac{x+y}{2a+1}.
\]
\end{theorem}
\begin{proof}
Poniendo \(X=\mathcal B^*\mathcal C+\mathcal C^*\mathcal B\), se obtiene
\[
 Q_-^*Q_-=(a^2+1)I-aX,\qquad
 Q_+^*Q_+=((a+1)^2+1)I+(a+1)X.
\]
Las ponderaciones cancelan \(X\) y el coeficiente restante es
\(2a^3+3a^2+3a+1=N\). La normalización da la isometría.
Sumar los dos canales recupera \((2a+1)\mathcal Bv\);
la combinación \(ay-(a+1)x\) recupera \((2a+1)\mathcal Cv\).
Finalmente, \(\mathcal B^*\mathcal B=I_H\).
\end{proof}

El balance vale también en dimensión infinita: sólo se componen
operadores acotados. La imagen de \(W_a\) es cerrada y su proyector
ortogonal es \(W_aW_a^*\). Por ello un par normalizado admite una
entrada exactamente cuando pertenece a esa imagen; su componente
ortogonal mide el defecto de compatibilidad.

La factorización
\[
 W_a=R_aE_a,\qquad
 E_av=\frac1{\sqrt p}
 \begin{pmatrix}\sqrt{a(a+1)}\,\mathcal Bv\\\mathcal Cv\end{pmatrix},
 \quad
 R_a=\frac1{\sqrt{2a+1}}
 \begin{pmatrix}\sqrt a&-\sqrt{a+1}\\
 \sqrt{a+1}&\sqrt a\end{pmatrix}
\]
se comprueba multiplicando los bloques; \(R_a\) actúa sobre
\(H'\oplus H'\), con cada entrada escalar multiplicando \(I_{H'}\).
En la partición nonádica
\(a=8\), los pesos son \(72/73\) y \(1/73\); el factor
\(73=8\cdot9+1\) es consecuencia de la ley general
\(a(a+1)+1\), no un dato primitivo añadido. El balance correspondiente es
\(9\|Q_-v\|^2+8\|Q_+v\|^2=17\cdot73\,\|v\|^2\).

\subsection{Transporte de observables y memoria entrante}
\label{iv:subsec:acoplamiento-lectores}

\begin{theorem}[Observables comunes a los dos canales]
\label{iv:thm:acoplamiento-lector}
Para un operador acotado autoadjunto \(F\) de \(H'\),
\[
 W_a^*\operatorname{diag}(F,F)W_a
 =\frac{a(a+1)\mathcal B^*F\mathcal B+\mathcal C^*F\mathcal C}{p}.
\]
Si \(\mathcal B\) es unitario y \(\mathcal C=\mathcal BR\), con
\(R\) unitario, y \(G=\mathcal B^*F\mathcal B\), el observable se
conserva si y sólo si \(GR=RG\).
\end{theorem}
\begin{proof}
Al expandir \((a+1)Q_-^*FQ_-+aQ_+^*FQ_+\), los términos cruzados
se cancelan. Los otros dos coeficientes son
\(a(a+1)(2a+1)\) y \(2a+1\). Dividir por \(N\) da la fórmula.
En el caso indicado, la igualdad con \(G\) equivale a \(R^*GR=G\);
la unitariedad la hace equivalente a \(GR=RG\).
\end{proof}

La norma se obtiene con \(F=I\). Un observable general requiere
la condición de conmutación: conservar norma no equivale a conservar
cada observable sin transformación.

\begin{theorem}[Extensión unitaria con memoria]
\label{iv:thm:acoplamiento-memoria}
Si \(H'=H\), \(\mathcal B=I\) y \(\mathcal C=U\) es unitario, sean
\[
 A=\frac{\sqrt{a+1}}{\sqrt N}Q_-,\qquad
 D=\frac{\sqrt a}{\sqrt N}Q_+,\qquad
 \mathscr R_a=\begin{pmatrix}A&-D^*\\D&A^*\end{pmatrix}.
\]
Entonces \(\mathscr R_a\) es unitario en \(H\oplus H\), su primera
columna es \(W_a\), y conserva \(\|v\|^2+\|m\|^2\) para toda
entrada \((v,m)\). La inversión es
\[
 v=A^*z_1+D^*z_2,\qquad m=-Dz_1+Az_2.
\]
\end{theorem}
\begin{proof}
Los operadores \(A,D\) son normales y conmutan entre sí y con sus
adjuntos, por ser polinomios de \(U,U^*\). El balance da
\(A^*A+D^*D=I\), y la normalidad da \(AA^*+DD^*=I\).
Los términos cruzados de ambos productos
\(\mathscr R_a^*\mathscr R_a\) y \(\mathscr R_a\mathscr R_a^*\)
se anulan; los diagonales son la identidad. Aplicar el adjunto a
\((z_1,z_2)\) proporciona la inversa.
\end{proof}

La recuperación del registro, la identificación orbital y la extensión
unitaria precisan tres niveles de información distintos. El primero
conserva los bloques con su carta; el segundo identifica operadores bajo
las hipótesis declaradas; el tercero transporta conjuntamente estado y
memoria. Esta jerarquía permite usar la misma estructura generada sin
confundir una lectura reducida con la totalidad del estado.
